\documentclass[10pt,a4paper]{amsart}
\usepackage[margin=19mm]{geometry}
\usepackage{amsmath,amssymb,mathtools}
\usepackage[scr=rsfs,cal=euler]{mathalfa}
\usepackage{xfrac}
\usepackage[unicode,hidelinks,bookmarks=false]{hyperref}
\newcommand{\ie}{i.e.\ }
\newcommand{\cf}{cf.\ }
\newcommand{\resp}{respectively\ }
\newcommand{\Subsection}{\subsection}
\providecommand{\nobreakdash}{\nobreak\hbox{-}}
\setlength{\emergencystretch}{2em}
\multlinegap0pt
\usepackage{bm}
\usepackage{bbm}
\usepackage{dsfont}
\usepackage{xspace}
\usepackage{cancel}
\usepackage{mathtools}
\usepackage{cleveref}
\usepackage{amsthm}
\makeatletter
\@ifundefined{proposition}{\newtheorem{proposition}{Proposition}}{}
\makeatother
\title[On the Picard number and the extension degree of period matr]{On the Picard number and the extension degree of period matrices of complex tori}
\author{Robert Auffarth, Jorge Duque Franco}
\date{Source: arXiv:2503.21642v2 (CC BY 4.0)}
\begin{document}
\maketitle

\noindent\emph{Source and licence.} On the Picard number and the extension degree of period matrices of complex tori. Version arXiv:2503.21642v2 (CC BY 4.0), distributed under the Creative Commons Attribution 4.0 International licence, as recorded in the arXiv licence metadata for this exact version and shown on the abstract page https://arxiv.org/abs/2503.21642v2. Full rights notice: licenses/arxiv-2503.21642-CC-BY-4.0.txt.\par\smallskip

\noindent\textit{Editorial scope.} Counted unit: the Proposition whose source label is \texttt{prop:2302} in the article (source0.tex lines 364--374, source label \texttt{prop:2302}), delivered together with the article's own complete original proof of it (source0.tex lines 376--394, one \texttt{proof} environment of 19 lines). No mathematics is written, repaired or re-derived: statement and proof are the article's own text line for line. Required context retained uncounted: none. Every definition, display and label of those lines is retained unchanged. Omitted from this package and not redistributed: 1--363, 375--375, 395--714. Nothing inside a retained excerpt is omitted or altered: every retained body is delivered exactly as the article's lines are written, so no omission receipt of any kind accompanies this package. Citations: the retained excerpts carry no citation command at all, so no bibliography entry of the article is transcribed and no reference is redistributed. Reference adaptation: none was needed and none was made; every reference inside the delivered unit resolves inside it, so no unresolved reference or citation is produced. Printed slips kept literal: no printed word, symbol, hypothesis or number of the counted statement or of its proof was altered. Layout and class: the article's own class is \texttt{article} with packages including graphicx, amsfonts, amsthm, amsxtra, amssymb, verbatim, makeidx, subeqnarray, relsize, euscript, babel, babelbib, inputenc, wrapfig; the wrapper is the lane's pinned \texttt{amsart} editorial wrapper at 10pt on A4, which restates the theorem-like environments the retained text needs and adds the article's own macro definitions that the retained text needs (none), with extra wrapper packages (bm, bbm, dsfont, xspace, cancel, mathtools, cleveref). The delivered numbering is the unit's own. Assets: no retained span contains a figure environment, a graphics-inclusion command, a PDF-page inclusion, a table or a TikZ picture; no image file is copied into this package, the package declares no asset dependency and no image file is redistributed with it. Licence: version arXiv:2503.21642v2 of this article is distributed under the Creative Commons Attribution 4.0 International licence, as recorded in the arXiv licence metadata for this exact version and shown on the abstract page https://arxiv.org/abs/2503.21642v2. A full rights notice is delivered at licenses/arxiv-2503.21642-CC-BY-4.0.txt. Characters: every retained excerpt is pure ASCII, so the delivered engine, XeLaTeX with its default Latin Modern text font, reports no missing character.
\begin{proposition}
\label{prop:2302}
Let $f:X\to Y$ be a map between complex tori with period matrices $(\tau\hspace{0.2cm}I_g)$ and $(\sigma\hspace{0.2cm}I_g)$, respectively. Then we have the following: 
\begin{enumerate}
\item If $f$ has finite kernel, then $F_X\subseteq F_Y$, and in particular $\mathfrak{d}_X\leq\mathfrak{d}_Y$.
\item If $f$ is surjective, then $F_Y\subseteq F_X$, and in particular $\mathfrak{d}_X\geq\mathfrak{d}_Y$.
\item If $f$ is an isogeny, then $F_X=F_Y$, and so $\mathfrak{d}_X=\mathfrak{d}_Y$.
\item $F_X=F_{X^\vee}$.
\end{enumerate}
In particular, 3. implies that $F_X$, and therefore $\mathfrak{d}_X$, only depends on $X$ and not on the specific period matrix chosen.
\end{proposition}

\begin{proof}
Clearly the third item follows from the previous two, and the fourth item is trivial since a period matrix for $X^\vee$ is simply $$(\tau^\tau\hspace{0.2cm}I_g).$$

To prove the first item, the homomorphism $f:X\to Y$ between complex tori induces the following equation:
\begin{equation}
\label{eq:0403}
    \rho_a(f)(\tau\hspace{0.2cm}I_g)=\left(\sigma\hspace{0.2cm}I_g\right)\rho_r(f),
    \end{equation}
where $\rho_a(f)$ is the analytic representation of $f$ and $\rho_r(f)$ is its rational representation. Writing 
$$\rho_r(f)=\left(\begin{array}{cc}A&B\\C&D\end{array}\right),\;\; A,B,C,D\in M_g(\mathbb{Z}),$$ 
we obtain in particular that
\[\rho_a(f)=\sigma B+D.\]
This implies that each coefficient of the matrix $\rho_a(f)$ belongs to the field $F_Y$. If $f$ has finite kernel, its analytic representation $\rho_a(f)$ is injective and therefore possesses a left inverse $M$ whose coefficients also lie in $F_Y$. Using \cref{eq:0403} we deduce that
\[\tau=M(\sigma A+C),\]
showing that each coefficient of $\tau$ also lies in $F_Y$. Therefore $F_X\subseteq F_Y$.

Now assume that $f$ is surjective. Then by dualizing, we get an injective morphism $f^\vee:Y^\vee\to X^\vee$, and so by the first and fourth items,
\[F_Y=F_{Y^\vee}\subseteq F_{X^\vee}=F_X.\]
\end{proof}

\end{document}
